\section{Exact uniform sampling}
\label{sec:exact-sampling}

We use the rare residual-mixture construction of G\"obel, Liu,
Manurangsi, and Pappik~\cite[Section~3.1, Algorithm~1 and
Theorem~6]{GLMP2024}. Their Corollary~15 computes the approximate law
by powers of the transition matrix. Exact tabulation is also used in
the contingency-table sampler~\cite[Sections~6.3--6.4]{OpenAIContingency2026}. Here we
establish the tabulation and bit-cost bounds directly for the rational
switch chain.

The idea is to sample from a very accurate finite-time law $p$ most of
the time, and to correct its remaining bias on a rare branch. For a
suitably small $\delta>0$, we will have $(1-\delta)p\le\pi_d$ pointwise.
Then
\[
 r=\frac{\pi_d-(1-\delta)p}{\delta}
\]
is a probability law, and the mixture $(1-\delta)p+\delta r$ is exactly
$\pi_d$. Only the branch of probability $\delta$ enumerates the state
space and tabulates $p$ to sample from $r$. We choose $\delta$
exponentially small enough to pay for that enumeration in expectation;
the chain still reaches the required accuracy in polynomially many
steps.

\begin{proof}[Proof of Corollary~\ref{cor:exact-sampling}]
We use the usual explicit binary encoding of the degree vector, with
delimiters between its entries. Let $n$ be the number of entries of
$d$. For $n<4$,
enumerate the at most $2^3$ simple graphs and retain those with the
specified labeled degrees. The unbiased integer draw described below
samples uniformly from this nonempty constant-size set. We may
therefore assume $n\ge4$ and use the notation of
Section~\ref{sec:introduction}. Havel's greedy degree-realization
procedure~\cite{Havel55}, with deterministic tie-breaking, constructs
one graph $G_0\in\Omega_d$ in polynomial bit time.

\paragraph{An accurate rational law.}
Set
\begin{equation}\label{exact:parameters}
\begin{gathered}
 B=\binom n2,\qquad M=|\Omega_d|\le2^B,\qquad
 q=12\binom n4,\\
 D=4B+10,\qquad \delta=2^{-D},\qquad
 t=\left\lceil\gamma_0^{-1}(B+2D+2)\right\rceil,
\end{gathered}
\end{equation}
where $\gamma_0=[24n^2\binom n4]^{-1}$ is from
Equation~\eqref{mix:gamma}. Thus $D=O(n^2)$ and $t=O(n^8)$.
Write $p_G=P_d^t(G_0,G)$. If $M>1$, Equation~\eqref{mix:TV} gives
\[
 \|p-\pi_d\|_{\mathrm{TV}}
 \le\frac12\sqrt M\,e^{-\gamma_0t}
 \le\frac12\,2^{B/2}e^{-(B+2D+2)}
 \le2^{-2D}=\delta^2.
\]
The last inequality follows from $2<e$ and $B,D\ge0$.
If $M=1$, then $p=\pi_d$, so the same bound holds.
Applying total variation to the singleton event $\{G\}$ gives
$p_G\le1/M+\delta^2$. Since $M\delta\le2^{B-D}\le1$,
\begin{equation}\label{exact:domination}
 (1-\delta)p_G
 \le\frac{1-\delta}{M}+(1-\delta)\delta^2
 \le\frac1M
 \qquad(G\in\Omega_d).
\end{equation}
The parameters in Equation~\eqref{exact:parameters}, apart from the
analytical quantity $M$, are computable from $n$ in polynomial bit
time. The ordinary sampling branch below does not need to know $M$.

For distinct $G,H$, the proof of Theorem~\ref{thm:main} shows that
$qP_d(G,H)$ is one for switch neighbors and zero otherwise. The row
sum identity gives
\[
 qP_d(G,G)=q-\bigl|\{H\ne G:H\text{ is a switch neighbor of }G\}\bigr|,
\]
which is also a nonnegative integer. Hence $A=qP_d$ is a
nonnegative integer matrix whose rows sum to $q$. Starting with the
integer point mass $a^{(0)}$ at $G_0$, define
\begin{equation}\label{exact:law}
 a^{(s+1)}=a^{(s)}A\quad(0\le s<t).
 \qquad
 p_G=\frac{a_G^{(t)}}{q^t},
 \qquad
 \sum_{G\in\Omega_d}a_G^{(t)}=q^t.
\end{equation}
All numerators are nonnegative integers.

We can tabulate this law exactly when needed. Enumerate all $2^B$
adjacency bit strings and retain exactly those whose labeled degrees
are $d$; this also computes $M$. Each degree test has polynomial bit
cost. For every retained pair of graphs, the switch-neighbor test and
the corresponding entry of $A$ are computable in polynomial bit time.
Dense propagation in Equation~\eqref{exact:law} uses at most $tM^2$
integer multiply-add operations. At step $s$, every numerator and
every partial nonnegative sum is at most $q^s$ or $q^{s+1}$,
respectively. Their bit lengths are therefore $O(t\log q)$, which is
polynomial in $n$. With, for example, schoolbook integer arithmetic,
the complete enumeration and tabulation cost at most
\begin{equation}\label{exact:tabulation-cost}
 2^{2B}\operatorname{poly}(n)
\end{equation}
bit operations. The enumerated graph universe depends only on $n$,
not on the accuracy parameter $D$.

\paragraph{Exact correction.}
We first specify the elementary random choice used throughout.
For an integer $K\ge1$, a uniform element of $\{0,\ldots,K-1\}$ is
obtained by drawing $\lceil\log_2 K\rceil$ independent unbiased bits
and repeating if the resulting integer is at least $K$; for $K=1$
use zero bits. Each attempt succeeds with probability greater than
$1/2$, except that a power of two succeeds with probability one.
The draw is exact, terminates almost surely, and uses fewer than two
attempts in expectation. An attempt has bit cost polynomial in
$\log(K+1)$.

Draw $D$ fresh unbiased bits. If they are not all zero, run exactly
$t$ steps of $P_d$ from $G_0$ and output the resulting graph.
Each step uses one bit for the holding decision and the preceding
integer draw to choose a uniform four-element vertex set and one of
the six ordered matching pairs. A four-set can be selected by
unranking its index in the lexicographic list, in polynomial bit
time. This implements $P_d$ exactly. The branch has output law $p$,
terminates almost surely, and has expected polynomial bit cost.

If the $D$ bits are all zero, perform the exact tabulation above and
write $a_G=a_G^{(t)}$. Form the integer weights
\begin{equation}\label{exact:weights}
 w_G=2^Dq^t-M(2^D-1)a_G,
 \qquad
 \sum_{G\in\Omega_d}w_G=Mq^t.
\end{equation}
They are nonnegative by Equation~\eqref{exact:domination}; their sum
uses $\sum_Ga_G=q^t$. Their bit lengths are
$O(D+t\log q+\log M)$, again polynomial in $n$. Draw a uniform
integer in $\{0,\ldots,Mq^t-1\}$ and select its interval in the
cumulative weights $w_G$. This samples exactly the probability law
\[
 r_G=\frac{w_G}{Mq^t}
     =\frac{1/M-(1-\delta)p_G}{\delta}.
\]
Use fresh random bits in all draws after the branch decision.
The correction branch has probability $\delta$, so the output law is
\begin{equation}\label{exact:mixture}
 (1-\delta)p_G+\delta r_G=\frac1M
 \qquad(G\in\Omega_d).
\end{equation}
Thus the output is exactly uniform.

\paragraph{Bit cost and termination.}
Every enumeration and integer computation on the correction branch
is finite. All rejection draws terminate almost surely; only finitely
many such draws are invoked on either branch. Hence the complete
algorithm terminates almost surely. The preprocessing, branch choice,
and ordinary branch have expected polynomial bit cost. By
Equation~\eqref{exact:tabulation-cost}, the conditional expected cost
of the correction branch, including its final exact integer draw and
cumulative scan, is at most $2^{2B}\operatorname{poly}(n)$. Its
contribution to the unconditional expectation is
\[
 \delta\,2^{2B}\operatorname{poly}(n)
 =2^{-2B-10}\operatorname{poly}(n)
 \le\operatorname{poly}(n).
\]
All loops, graph tests, and integer computations are specified by one
algorithm, and the input explicitly lists $n$ degrees. The resulting
expected bit bound is therefore a fixed polynomial in the input
length. The correction branch may take exponential time, and the
rejection draws have unbounded tails; the conclusion is an expected
bound, not a polynomial bound on every execution.
\end{proof}
